% Part: incompleteness
% Chapter: theories-computability
% Section: complete-decidable

\documentclass[../../../include/open-logic-section]{subfiles}

\begin{document}

\olfileid{inc}{tcp}{cdc}

\olsection{Teori kang \printtoken{S}{axiomatizable} lan jangkep bisa diputusake}

Teori diarani {\em jangkep} yen kanggo saben ukara $!A$,
$!A$ utawa $\lnot !A$ bisa dibuktekake.

\begin{lem}
Upamana teori $\Th{T}$ jangkep lan !!{axiomatizable}. Mula
$\Th{T}$ bisa diputusake.
\end{lem}

\begin{proof}
Upamana $\Th{T}$ jangkep lan $A$ himpunan aksioma kang komputabel.
Yen $\Th{T}$ ora konsisten, cetha yen iki komputabel. (Algoritmane:
``cukup wangsulana ya.'') Mula kita bisa nganggep yen $\Th{T}$
uga konsisten.

Kanggo mutusake apa ukara $!A$ kalebu ing~$\Th{T}$ utawa ora,
goleka bebarengan !!a{derivation} kanggo~$!A$ saka~$\Th{T}$ lan
!!a{derivation} kanggo~$\lnot !A$. Amarga $\Th{T}$ jangkep, mesthi
nemokake salah sijine; lan amarga $\Th{T}$ konsisten, yen sampeyan
nemokake !!a{derivation} kanggo~$\lnot !A$, ora ana !!{derivation} kanggo~$!A$.

Kanthi tembung liya, kita wis ngerti yen $\Th{T}$ iku !!{c.e.}; mula,
miturut teorema kang wis dibuktekake sadurunge, cukup dituduhake yen
komplemen saka~$\Th{T}$ uga !!{c.e.}. Nanging !!a{formula}~$!A$ kalebu
ing~$\Th{\bar T}$ yen lan mung yen $\lnot !A$ kalebu ing~$\Th{T}$; mula
$\Th{\bar T} \leq_m \Th{T}$.
\end{proof}

\end{document}
